\begin{figure}[!htb]
\centering
\begin{tikzpicture}[x=1mm,y=1mm,font=\sffamily\small]
 \fill[paleblue] (8,6)--(56,6)--(68,42)--(20,42)--cycle;
 \draw[teal,line width=.7pt] (8,6)--(56,6)--(68,42)--(20,42)--cycle;
 \draw[->,>=stealth,muted] (8,6)--(84,6) node[right] {$x$};
 \draw[->,>=stealth,muted] (8,6)--(8,49) node[above] {$y$};
 \draw[->,>=stealth,brand,line width=1pt] (8,6)--(32,6);
 \draw[->,>=stealth,brand,line width=1pt] (32,6)--(38,24);
 \draw[->,>=stealth,teal,line width=.8pt] (8,6)--(38,24);
 \draw[dashed,muted] (38,24)--(38,6);
 \fill[brand] (38,24) circle (1.2);
 \node[anchor=north,text=brand] at (22,5) {$u\mathbf a$};
 \node[anchor=west,text=brand] at (38,14) {$v\mathbf b$};
 \node[anchor=south east,text=teal] at (25,16) {$\mathbf r$};
 \node[anchor=north east] at (7,5) {$O$};
 \node[anchor=west,text=brand] at (42,27) {$\mathbf f=(u,v,0)$};
 \node[anchor=west,text=muted,font=\sffamily\footnotesize] at (78,35) {晶胞在平面内的截面};
 \node[anchor=west,text=muted,font=\sffamily\footnotesize] at (78,28) {箭头相加得到位置};
 \node[anchor=north,text=teal] at (54,4) {$\mathbf a$};
 \node[anchor=east,text=teal] at (15,31) {$\mathbf b$};
\end{tikzpicture}
\caption{分数坐标与笛卡尔位置的关系（二维示意）。$u\mathbf a$ 与 $v\mathbf b$ 相加得到 $\mathbf r$；斜晶胞的 $\mathbf b$ 方向同时含有全局 $x$ 和 $y$ 分量。}
\label{fig:round10-fractional-cartesian}
\end{figure}
